\section{第二次创业：AI 人口、Will 与 Skill}

产品筛选标准讲完后，问题进入更宏观的判断：这一代 AI 到底在创造什么。张月光的答案不是“更多服务”，而是“AI 人口”。

张月光第二次创业时，有一个更宏大的判断：这代 AI 技术不是创造 AI 服务，而是创造 AI 人口。所谓 AI 人口，是大量可交互、可执行、可协作的 AI 个体进入人类网络。这个判断解释了为什么他会同时做 AI 乙女游戏和 Agent：前者探索情感/内容里的 AI 个体，后者探索生产力里的 AI 个体。

\begin{figure}[H]
\centering
\includegraphics[width=0.94\textwidth]{ai-population-will-skill.png}
\caption{AI 人口、Will 与 Skill：人负责目标，AI 负责执行。}
\end{figure}

\begin{knowledgebox}{读图：AI 服务和 AI 人口的差别}
如果 AI 是服务，用户调用功能；如果 AI 是人口，用户和某个会做事的对象协作。图中 Will 指目标、意图、选择和品味；Skill 指执行、生成、工具调用和越过个人能力边界。
\end{knowledgebox}

\subsection{人负责 Will，AI 负责 Skill}

上一节把 AI 看成“人口”，这一节拆开人与 AI 的分工。Will / Skill 是一个简短但很实用的产品框架。

访谈中他说，人应该负责 Will，AI 负责 Skill。这句话很适合解释 Agent 产品：用户不用掌握所有技能，但要知道自己想要什么、什么结果好、什么时候继续、什么时候停止。AI 则承担执行、工具调用和生成。

\begin{importantbox}{Will / Skill 分工}
AI 产品越强，越不应把人完全排除在外。人的价值转向目标、判断、品味和选择；AI 的价值转向执行、生成和能力扩展。
\end{importantbox}

\subsection{本章小结}

AI 人口不是一句口号，而是产品定位变化：用户不只是使用功能，而是与 AI 个体协作。Will / Skill 分工，是张月光理解 Agent 和 AI 朋友的核心框架。

